% ==================================================================
% id:       func-026
% topic:    מספרים ממשיים
% subtopic: מספרים רציונליים ואי־רציונליים
% solution: full
% ==================================================================
\textbf{יהיו
$q_1,q_2\in\mathbb{Q}$
ויהיו
$r_1,r_2\notin\mathbb{Q}$.}
\begin{enumerate}
\item
הראו כי
$q_1+q_2\in\mathbb{Q}$,
כי
$q_1q_2\in\mathbb{Q}$,
וגם כי
$\frac{q_1}{q_2}\in\mathbb{Q}$
אם
$q_2\neq0$.
\hint{כתבו כל מספר רציונלי כמנה של שני שלמים, ובדקו מה מתקבל בכל פעולה.}
\begin{solbox}
נתון כי
$q_1,q_2\in\mathbb{Q}$
ולכן קיימים
$m_1,n_1,m_2,n_2\in\mathbb{Z}$
עבורם
$n_1,n_2\neq0$
ו-
$$.q_1=\frac{m_1}{n_1}\quad,\quad q_2=\frac{m_2}{n_2}$$
מכאן נחבר ונכפיל ונקבל
$$q_1+q_2=\frac{m_1n_2+m_2n_1}{n_1n_2}\ ,\ q_1q_2=\frac{m_1m_2}{n_1n_2}.$$
המונים והמכנים הם מספרים שלמים, והמכנה
$n_1n_2\neq0$,
ולכן שני המספרים רציונליים.

לחילוק: אם
$q_2\neq0$
אז
$m_2\neq0$,
ולכן
$$\frac{q_1}{q_2}=\frac{m_1}{n_1}\cdot\frac{n_2}{m_2}=\frac{m_1n_2}{n_1m_2},$$
כאשר
$m_1n_2$
ו-$n_1m_2$
שלמים ו-$n_1m_2\neq0$.
לכן גם
$\frac{q_1}{q_2}$
רציונלי.
\end{solbox}
\item
הראו כי
$q_1+r_1\notin\mathbb{Q}$,
ושאם
$q_1\neq0$
אז גם
$\frac{r_1}{q_1}\notin\mathbb{Q}$.
\hint{הניחו בשלילה שהתוצאה רציונלית, והשתמשו בסעיף הקודם.}
\begin{solbox}
נניח בשלילה כי
$q_1+r_1\in\mathbb{Q}$,
וניזכר כי
$q_1\in\mathbb{Q},\ r_1\notin\mathbb{Q}$.
מכאן נסיק כי חיסור רציונליים יישאר רציונלי, כלומר
$$r_1=(q_1+r_1)-q_1\in\mathbb{Q}$$
בסתירה להנחה שזהו מספר אי-רציונלי.
באופן דומה, אילו
$\frac{r_1}{q_1}\in\mathbb{Q}$
וגם
$q_1\neq0$
אז נקבל מסגירות הרציונליים לכפל שבהכרח
$r_1=q_1\cdot\frac{r_1}{q_1}\in\mathbb{Q}$
וזוהי סתירה.
\end{solbox}
\item
הראו כי לא בהכרח מתקיים
$r_1+r_2\notin\mathbb{Q}$.
\hint{חפשו שני מספרים אי-רציונליים שהחלק האי-רציונלי שלהם מתבטל בסכום.}
\begin{solbox}
חיבור של מספרים אי-רציונליים לא בהכרח ייצא מספר אי-רציונלי. למשל,
$r_1=\sqrt2,\ r_2=-\sqrt2$
שניהם אי-רציונליים, אך
$r_1+r_2=0$
וזהו מספר רציונלי.
\end{solbox}
\end{enumerate}
